\section{System Overview}\label{sec:overview}

\Cref{fig:overview-loop} shows one v4 cycle. Generation $k$ is the current champion,
stored as immutable ONNX, PyTorch and checkpoint files. It plays self-play games in the Go
engine, which writes one compressed shard per cycle (\cref{sec:selfplay}). The learner
warm-starts from generation $k$'s weights and trains on a window over all shards
(\cref{sec:learner}). It then exports a candidate, and that export must pass a numeric
parity check. The orchestrator gates the candidate against generation $k$ with a
sequential test. If it passes, the candidate becomes generation $k{+}1$ and is published
to a champion registry (\cref{sec:orchestration}). Meanwhile, self-play for the next cycle
is already running.

\begin{figure}[htbp]
\centering
\begin{tikzpicture}[
  node distance=5mm and 7mm,
  box/.style={draw, rounded corners=2pt, align=center, font=\small, minimum height=9mm, inner sep=3pt},
  data/.style={draw, align=center, font=\small, minimum height=7mm, inner sep=2pt, fill=gray!10},
  arr/.style={-{Stealth[length=2mm]}, thick},
]
\node[box] (sp) {Self-play\\\footnotesize Go, PUCT MCTS\\\footnotesize cap randomization};
\node[data, right=of sp] (sh) {shard\\\footnotesize\code{cycle-N.npz}};
\node[box, right=of sh] (tr) {Learner\\\footnotesize masked loss, D4,\\\footnotesize EMA, window};
\node[box, right=of tr] (ex) {Export\\\footnotesize ONNX + parity};
\node[box, below=9mm of ex] (ga) {Gate\\\footnotesize SPRT vs.\ gen $k$};
\node[box, left=of ga] (pr) {Promote\\\footnotesize gen $k{+}1$};
\node[box, left=of pr] (pu) {Publish\\\footnotesize registry, release,\\\footnotesize regression check};
\node[data, below=9mm of sp] (ch) {champion\\\footnotesize gen $k$};
\draw[arr] (sp) -- (sh);
\draw[arr] (sh) -- (tr);
\draw[arr] (tr) -- (ex);
\draw[arr] (ex) -- (ga);
\draw[arr] (ga) -- node[above, font=\footnotesize]{accept} (pr);
\draw[arr] (pr) -- (pu);
\draw[arr] (pu) -- (ch);
\draw[arr] (ch) -- node[left, font=\footnotesize]{evaluator} (sp);
\draw[arr, dashed] (ch.east) to[out=0, in=250] node[below right, font=\footnotesize, pos=0.35]{\code{--init-from}} (tr.south);
\end{tikzpicture}
\caption{One v4 training cycle. Every arrow is a resumable stage recorded in the run's
state file. Self-play for cycle $N{+}1$ starts while cycle $N$ trains and gates.}
\label{fig:overview-loop}
\end{figure}

\paragraph{Interfaces.}
The three components meet at two narrow contracts. That let two developers (here, two
agent sessions) build them in parallel and test them separately.
\begin{enumerate}
  \item \textbf{Data.} GOFER SHARD v1 is a \code{.npz} archive with a fixed set of typed
        arrays (\cref{tab:selfplay-shard} in \cref{sec:selfplay}). The Go engine writes
        it, and the learner's packer can also produce it from legacy JSONL.
  \item \textbf{Learner command line.} The trainer takes \code{--data} (files or directories),
        \code{--window-rows}, \code{--init-from} or \code{--continue}, and an optional
        \code{--config} preset. It writes \code{best.pt}, a self-describing
        \code{best.ckpt}, a full-state \code{last.pt}, and a \code{summary.json} with
        fixed keys. The exporter takes a checkpoint and writes an ONNX file whose input and
        output names match the engine's schema, exiting non-zero if parity fails.
\end{enumerate}

\paragraph{Comparison.}
\Cref{tab:overview-compare} places v4 next to the v3 loop and the published \katago{}
recipe~\citep{wu2019accelerating}, one technique per row. Four symbols are used:
\checkmark{} means adopted as published, $\sim$ means adopted with a stated change,
\texttimes{} means absent, and ``new'' means not in the \katago{} paper.

\begin{table}[htbp]
\centering
\footnotesize
\begin{tabular}{@{}>{\raggedright\arraybackslash}p{3.9cm}>{\raggedright\arraybackslash}p{3.0cm}>{\raggedright\arraybackslash}p{2.6cm}>{\raggedright\arraybackslash}p{5.5cm}@{}}
\toprule
Technique & \katago{} paper & \gofer{} v3 & \gofer{} v4 \\
\midrule
Playout cap randomization & $p{=}0.25$, fast turns discarded & $p{=}0.20$, discarded & $\sim$ $p{=}0.20$; fast turns kept, policy loss masked (\S\ref{sec:learner-masking}) \\
Forced playouts + target pruning & \checkmark & $\sim$ simplified & $\sim$ unchanged: forced root visits; children with ${<}2$ visits dropped from target (\S\ref{sec:selfplay-search}) \\
Global pooling bias & \checkmark{} (mean, scaled mean, max) & \texttimes & $\sim$ mean + max; \code{gpool} nets only \\
Auxiliary opponent policy & \checkmark{} $w{=}0.15$ & \texttimes{} (label wrong) & \checkmark{} $w{=}0.15$, \code{gpool} nets \\
Ownership target & \checkmark{} cross-entropy & $\sim$ MSE, absolute frame & $\sim$ MSE on $\tanh$, side-to-move frame \\
Score target & pdf + cdf losses & \texttimes & $\sim$ Huber on margin$/20$ \\
Go-specific input features & liberties, ladders, \dots & \texttimes & \texttimes{} (schema v2 unchanged) \\
Rules / board-size randomization & \checkmark & \texttimes & \texttimes{} (net is size-agnostic) \\
Symmetry augmentation & --- & \texttimes & new: per-sample D4 \\
Replay window & power law, eq.~\eqref{eq:window} & FIFO 50k rows & \checkmark{} power law + recency weighting \\
Validation split & --- & by row (leaks) & new: by game \\
Weight averaging & EMA of snapshots & \texttimes & $\sim$ per-step EMA, $0.999$ \\
LR schedule & fixed, reduced early/late & constant & $\sim$ warmup + cosine per cycle \\
Architecture change & train next size in parallel & re-bootstrap & new: distillation from champion \\
Gating & ${\ge}100/200$ wins & 0.55 + Wilson, ${\le}200$ games, interim stops & new: SPRT + Wilson fallback; fewer false promotions \\
Preemption safety & --- & restart cycle & new: stage + step resume \\
Data format & --- & JSONL & new: binary shards, ${\sim}40\times$ smaller \\
\bottomrule
\end{tabular}
\caption{Technique-by-technique comparison. ``---'' means the \katago{} paper does not
describe the item. The v4 column covers only what is implemented and tested. None of the
v4 differences has been tested for playing strength (\cref{sec:limitations}).}
\label{tab:overview-compare}
\end{table}

The table is deliberately balanced. On the core learning techniques, v4 is a
\emph{subset} of \katago{} with a few simplifications. Policy target pruning is only a
visit-threshold approximation, and the score-belief distribution, Go-specific features
and rules randomization are all missing. The additions
are mostly engineering: data format, resumability, sequential gating, distillation. These
matter when compute is scarce and preemptible. They do not change what the network can
learn in principle.
